\scanpage{84}
\mspage{18}
On a alors (cf (56) pour la définition de $\rho_s$)
\[
(\unsure{76|75})\qquad
\rho_s=\varepsilon_s^2\deleted{=\rho_s}\inserted{=\sigma_a\rho\sigma_a\rho}
\quad\deleted{(\text{cf }\unsure{56})}
\]
\margin{\uline{NB}
$\rho_0=\sigma_{a_0}\rho\sigma_{a_0}\rho$
\quad
$\rho_1=\rho(\rho_0)=\rho\sigma_{a_0}\rho\sigma_{a_0}$}
Cela \inserted{((76) et (73 c))} montre que $\rho,\sigma_{a_0}$ engendrent $\Pi_P^D$, et on trouve facilement
\[
\textcircled{77}\qquad
\Pi_P^D\simeq(\rho,\sigma_{a_0}\mid\rho^n=\sigma_{a_0}^2=1)
\simeq\mathbb Z/n\mathbb Z*\mathbb Z/2\mathbb Z
\]
\margin{\uline{NB} Pour $n=3$, on identifie \unsure{dans la notation ancienne de $P$, $\infty$ au} sommet origine, et $\sigma_{a_0}$ se note alors $\sigma_\infty$.}
Car la relation (59) est déjà conséquence, moyennant (76) et (73 c), des deux relations fondamentales
\[
(78)\qquad \rho^n=1,\quad\sigma_{a_0}^2=1.
\]
Les formules (73) se complètent en
\[
(79)\qquad
\left\{
\begin{aligned}
\rho(\varepsilon_s)&=\varepsilon_{\rho(s)}\\
\rho(\rho_s)&=\rho_{\rho(s)}
\end{aligned}
\right.
\qquad(s\in S)
\]
Conséquence immédiate de (73)(74)(75), ce qui nous montre la façon dont $\rho$ opère sur $\Pi_P$.

\deleted{Quant aux $\widetilde\sigma_b$ \inserted{($b\in A$)}, ils opèrent par transport des structures
\[
(80)\qquad
\left\{
\begin{aligned}
\widetilde\sigma_b(\varepsilon_s)&=\varepsilon_{\sigma_b^0(s)}^{-1}\\
\widetilde\sigma_b(\rho_s)&=\rho_{\sigma_b^0(s)}^{-1}
\end{aligned}
\right.
\qquad\forall\,s\in S,\ b\in A
\]
Il reste à donner l'opération}
On trouve (conséquence de (65) et (62 bis))
\[
(80)\quad
\left\{
\begin{aligned}
\tau_{a_0}(\rho_0)&=\rho_0^{-1},\quad
\tau_{a_0}(\rho_1)=\rho_1^{-1},\quad
\tau_{a_0}(\rho_2)=\rho_1^{-1}(\rho_2^{-1})\\
\cdots\quad\tau_{a_0}(\rho_i)&=
(\rho_{i-1}\rho_{i-2}\cdots\rho_1)^{-1}(\rho_i^{-1})\\
&=\underbrace{(\rho_{n\deleted{\unsure{-1}}}\cdots\rho_{i+1})}_{\text{\unsure{commencer} avec }\rho_n=\rho_0,\ \uline{\text{pas}}\ \rho_{n-1}!}(\rho_i^{-1})
\end{aligned}
\right.
\quad 0\leq i\leq n-1
\]
\margin{À retenir :
$\tau_a(\rho_s)=\rho_s^{-1}$ si $s\in S$ incident à $a\in A$}
\margin{Attention, je préfère écrire $\tau_{a_0}$ au lieu de $\tau_0$, (car pour $n=3$ la ``notation reçue'' pour $\tau_{a_0}$ est $\tau_\infty$, et non $\tau_0$ !)}
\margin{\uline{NB} L'action de $\tau_a$ sur $\Pi_P$ est la composée de
\[
\tau:\Pi_1(\Sigma^*,P_\omega)\longrightarrow\Pi_1(\Sigma^*,P_{\omega'})
\]
\unsure{avec}
\[
\lambda_a:\Pi_1(\Sigma^*,P_{\omega'})\longrightarrow\Pi_1(\Sigma^*,P_\omega)
\]}
\deleted{Pour} Les relations (79) se généralisent aussitôt en
\[
\deleted{(\unsure{81|80})}\qquad
\left\{
\begin{aligned}
\widetilde g(\varepsilon_s)&=(\varepsilon_{\widetilde g^0(s)})^{\operatorname{sg}(g)}\\
\widetilde g(\rho_s)&=(\rho_{\widetilde g^0(s)})^{\operatorname{sg}(g)}
\end{aligned}
\right.
\qquad
\begin{gathered}
(g\in D_P)\\
s\in S
\end{gathered}
\]
\deleted{$\unsure{\widetilde\sigma_a}$}
\margin{\uline{NB} L'action de $\widetilde D_P\subset\Pi_P^{D,\tau}$ sur $\Pi$ est par ``transport de structures'', i.e. celle de $\widetilde D_P^0\subset D_P\times\{1,\tau\}$ sur $(\Sigma^*,P_\omega)$.}
\margin{83}
